\authorsnames[{1,2,3,4,5},1]{Jerold Pearson, Roger E. Levine, Jon A. Krosnick, Resty Fufunan, Josearmando Torres}

\authorsaffiliations{
  {Stanford Alumni Association, Stanford University, Stanford, CA, U.S.A.},
  {American Institutes for Research, Palo Alto, CA, U.S.A.},
  {Department of Communication, Stanford University, Stanford, CA, U.S.A.},
  {Stanford Institute for Research in the Social Sciences, Stanford University, Stanford, CA, U.S.A.},
  {Stanford Institute for Research in the Social Sciences, Stanford University, Stanford, CA, U.S.A.}
}

\leftheader{Pearson, Levine, Krosnick, Fufunan, and Torres}
\rightheader{The Impact of Lottery Incentives}

\title{The Impact of Lottery Incentives: Increasing Response Rates Can Decrease Survey Accuracy}

\authornote{Contact information:
Jerold Pearson, Stanford Alumni Association, Stanford University, Stanford, CA 94305 (E-Mail: jpearson@alumni.stanford.edu).
Roger E.\ Levine, American Institutes for Research, 1070 Arastradero Road, Suite 200, Palo Alto, CA 94304 (E-Mail: x).
Jon A.\ Krosnick, Department of Communication, Stanford University, Stanford, CA 94305 (E-Mail: krosnick@stanford.edu).
Resty Fufunan, Stanford Institute for Research in the Social Sciences, Stanford University, Stanford, CA 94305 (E-Mail: resty@stanford.edu).
Josearmando Torres, Stanford Institute for Research in the Social Sciences, Stanford University, Stanford, CA 94305 (E-Mail: jtorres0@stanford.edu).}

\abstract{When offering incentives with a fixed budget, survey researchers can conduct
lotteries with many winners of smaller prizes or with fewer winners of larger prizes.
The expected value of the reward for a respondent is the same, but one approach may
be more psychologically compelling than the other. In an online survey of a listed
sample, the present study randomly assigned each potential respondent to be offered
five chances to win \$100, one chance to win \$500, or no incentive. Lottery incentives
produced a marginal increase in overall response rates (32.5\% vs.\ 30.6\%,
$p = 0.055$), and the two prize structures performed identically. Incentivized
respondents completed the survey more thoroughly and their responses were statistically
indistinguishable from those of non-incentivized respondents. More importantly,
incentives disproportionately increased participation among subgroups already
over-represented in the no-incentive condition---particularly women and donors---widening
rather than narrowing existing participation gaps. Researchers should assess baseline
heterogeneity in response propensity before deploying lottery incentives, as such
strategies may increase overall response while simultaneously amplifying nonresponse bias.}

\keywords{survey incentives; lottery; response rates; nonresponse bias; sample composition}


%%% Local Variables:
%%% mode: latex
%%% TeX-master: "srm_main"
%%% End:
